\documentclass[12pt,a4paper]{article}
\usepackage[T1]{fontenc}
\usepackage[utf8]{inputenc}
\usepackage[brazilian,english]{babel}
\usepackage{amsmath,booktabs,geometry,setspace,array,hyperref}
\geometry{top=3cm,bottom=2cm,left=3cm,right=2cm}
\onehalfspacing
\setlength{\parindent}{1.25cm}
\hypersetup{colorlinks=true,linkcolor=black,urlcolor=blue,citecolor=black}
\title{Práticas abertas para a documentação de pesquisas}
\author{João Silva\thanks{joao.silva@example.org} \and Ana Souza}
\date{}

\begin{document}
\begin{center}
  {\large\bfseries INSTITUTO DE EXEMPLO}\par
  \vspace{6pt}
  {\Large\bfseries \MakeUppercase{Práticas abertas para a documentação de pesquisas}}\par
  \vspace{6pt}
  João Silva\textsuperscript{1} \quad Ana Souza\textsuperscript{2}
\end{center}
\noindent\textsuperscript{1}Departamento de Métodos, Universidade A.\\
\textsuperscript{2}Laboratório de Dados, Universidade B.

\begin{abstract}
Este artigo apresenta um roteiro para registrar decisões, transformações e
resultados de uma pesquisa computacional. O modelo combina metadados,
descrição metodológica e uma tabela de rastreabilidade em uma estrutura
compatível com submissões acadêmicas brasileiras.
\end{abstract}
\noindent\textbf{Palavras-chave:} documentação; ciência aberta; métodos.

\selectlanguage{english}
\begin{abstract}
This article presents a workflow for recording decisions, transformations, and
results in computational research. The model combines metadata, a methods
description, and a traceability table in a structure suitable for academic
submissions.
\end{abstract}
\noindent\textbf{Keywords:} documentation; open science; methods.
\selectlanguage{brazilian}

\section{Introdução}
Uma pesquisa pode ser tecnicamente correta e ainda assim difícil de revisar
quando as decisões intermediárias não são registradas. Este modelo separa a
descrição do problema, os procedimentos e as evidências que sustentam a
conclusão.

\section{Procedimentos metodológicos}
Para cada observação $x_i$, o valor normalizado é definido por
\begin{equation}
  z_i = \frac{x_i-\mu}{\sigma},
  \label{eq:normalization}
\end{equation}
onde $\mu$ e $\sigma$ devem ser calculados apenas no conjunto de treinamento.
Essa convenção torna o exemplo reproduzível e explicita uma decisão que
frequentemente fica escondida no texto.

\begin{quote}
  Citações diretas longas devem ser apresentadas em bloco, com a fonte e a
  referência completas, conforme as normas adotadas pela revista.
\end{quote}

\section{Resultados e discussão}
A Tabela~\ref{tab:trace} registra as etapas que levam aos resultados. Em um
artigo real, substitua os valores ilustrativos pelos dados da pesquisa.
\begin{table}[ht]
  \centering
  \caption{Registro de rastreabilidade do estudo.}
  \label{tab:trace}
  \begin{tabular}{p{0.25\linewidth}p{0.27\linewidth}p{0.30\linewidth}}
    \toprule
    Etapa & Entrada & Evidência produzida \\
    \midrule
    Limpeza & 240 observações & conjunto validado \\
    Modelagem & conjunto validado & estimativas e erros \\
    Revisão & estimativas & relatório e dados \\
    \bottomrule
  \end{tabular}
\end{table}

\section{Considerações finais}
O autor deve conferir a paginação, as chamadas de figuras e tabelas, os
metadados e a versão vigente das normas antes da submissão. A equação
\eqref{eq:normalization} exemplifica como manter as escolhas analíticas
visíveis no corpo do texto.

\section*{Referências}
\begin{enumerate}
  \item J. W. Creswell, \emph{Research Design}, Sage, 2018.
  \item R. D. Peng, ``Reproducible research in computational science,''
  \emph{Science}, 2011.
\end{enumerate}
\end{document}
